\section{Related work}
\label{sec:related}
\subsection{Transactions, update provenance, and event histories}
Transactions provide atomic, durable state transformations \citep{gray1981transaction}, and optimistic concurrency control validates conflicting work before effects are installed \citep{kung1981occ}. These mechanisms support complete successors and commit-time predecessor checks. They do not determine the object of human authorization: an internally consistent transaction can admit a context-equivalent candidate that was never the reviewed instance.

Database provenance explains contributing data, derivations, and source locations \citep{buneman2001whywhere,cheney2009provenance}. Curated-database provenance also records edits and copies across versions; Buneman et al. explicitly connect unchanged data to the preceding version and support source, history, and modification queries \citep{buneman2006curated}. Reenactment reconstructs multi-version transaction provenance \citep{arab2018reenactment}, while PROV-DM represents entities, activities, agents, derivation, and responsibility \citep{moreau2013provdm}. Exact source preservation and rich lineage are established capabilities. The question here is how they participate in an admission condition that preserves the machine proposal while assigning authority to a possibly different human-authorized value and the exact reviewed candidate.

Event sourcing reconstructs state from events \citep{fowler2005eventsourcing}. An event journal can preserve proposals, corrections, reviews, versions, and field sources. Its persistence representation does not settle whether approval targets equivalent content or a particular instance. E1 holds the journal design and transactional mechanisms constant while adding the reviewed-candidate link. Agreement with the relational reference shows that candidate-bound authorization can be enforced in both representations; equal-valued substitution isolates the policy difference.

\subsection{Human approval, correction, and data curation}
Human-controlled integrity transitions have substantial precedents. Clark--Wilson combines well-formed transactions, authenticated execution, separation of duty, and reconstructible audit records \citep{clark1987integrity}. Mershad et al. persist cell versions pending principal-investigator approval; selection by cell identifier and timestamp already supports version-specific content approval \citep{mershad2015approval}. The present distinction concerns the authorization target when multiple persisted proposals share both values and retained review context.

Interactive cleaning separates system suggestions from human judgments. Guided Data Repair lets users confirm or reject repairs, retain a value, or supply a replacement \citep{yakout2011guided}. Falcon converts corrections into candidate rules and recognizes that identical current effects can conceal different semantic validity \citep{he2016falcon}. Wrangler provides editable transformation histories, previews, and reusable scripts \citep{kandel2011wrangler}; DBWiki combines identities, version history, atomic edits, and provenance queries \citep{buneman2011dbwiki}. These works establish correction and curation capabilities. We contribute an admission relation and its information characterization: preserve the original proposal, bind the explicitly authorized value to the exact reviewed instance, and carry both roles into a fresh complete successor with total field sources. The paired histories explain what disappears when these distinctions are omitted.

\subsection{Governed AI state and persistent memory}
Continuity Kernel is a close antecedent: it binds exact proposal identity, predecessor, pre-state authority, signed approver evidence, and lineage to an atomic complete accepted unit, with human preparation and multiple storage realizations \citep{he2026continuity}. Exact proposal binding is therefore an established neighboring principle. Our contribution specifies the correction-aware field relation and characterizes its failure-distinguishing information: proposed and authorized values retain separate roles, changed and unchanged fields retain exact sources, and equal-valued substitution isolates the required authorization target.

MemTxn separates admission, temporal visibility, and complete-state recovery \citep{cui2026memtxn}. MutMem binds signed retrieval-weight mutations to memory identity, provenance, signer epochs, and predecessor continuity \citep{saidi2026mutmem}. LatticeMind manages claim status, conflict, supersession, and provenance \citep{zhou2026latticemind}; Stored Is Not Supported governs assertion release through typed provenance from an authenticated accepted head \citep{he2026stored}. These preprints address adjacent stages of persistent-state governance and may encode the relation studied here. Our claim concerns its explicit correction semantics, information characterization, and executable separation of content/context approval from instance authorization.

\subsection{Binding and executable specification}
DPoP binds possession proofs to request method and URI \citep{fett2023dpop}; HTTP Message Signatures cover application-selected message components \citep{backman2024signatures}. OWASP transaction-authorization guidance requires server-controlled approval data, invalidation after relevant changes, and a final gate \citep{owaspTransactionAuthorization}. E2 applies this binding discipline to transfer an exact candidate review into confirmation. Its request and display controls distinguish server-record consistency from faithful presentation to a reviewer.

Alloy and Kodkod support relational specification and bounded model finding \citep{jackson2002alloy,torlak2007kodkod}, while TestEra connects specifications to execution through concretization and abstraction \citep{marinov2001testera}. We use these methods to examine the admission relation. The paired-history characterization provides the information distinction; the formal and persisted-state checks connect it to executable enforcement. Formal scope and implementation correspondence are stated in Sections~3, 4, and 6.
